\section{Aplicaciones Modernas}

\begin{frame}{Primeros Desarrollos y Aplicaciones (1950 - 1970)}
  El periodo de posguerra requirió materiales con un rendimiento superior al de
  los metales en condiciones específicas.

  \begin{block}{Material: Alúmina (\ce{Al2O3})}
    \begin{itemize}[<+->]
      \item \textbf{Propiedad Clave:} Alto aislamiento dieléctrico y
        resistencia térmica.
      \item \textbf{Aplicación:} Aisladores de bujías en motores y sustratos
        para circuitos electrónicos.
    \end{itemize}
  \end{block}

  \begin{block}{Material: Circonia (\ce{ZrO2})}
    \begin{itemize}[<+->]
      \item \textbf{Propiedad Clave:} Muy baja conductividad térmica y alta
        estabilidad.
      \item \textbf{Aplicación:} Recubrimientos de barrera térmica (TBC) en
        componentes de motores a reacción.
    \end{itemize}
  \end{block}
\end{frame}

\begin{frame}{Aplicaciones en Electrónica y Biomedicina (1980 - 2000)}
  \begin{block}{Componentes Electrónicos y Sensores}
    \begin{itemize}[<+->]
      \item \textbf{Materiales:} Titanato de Bario (\ce{BaTiO3}),
        Titanato-zirconato de plomo (PZT).
      \item \textbf{Propiedades Clave:} Alta constante dieléctrica y efecto
        piezoeléctrico.
      \item \textbf{Aplicaciones:} Condensadores cerámicos multicapa (MLCC),
        transductores para ultrasonido.
    \end{itemize}
  \end{block}

  \begin{block}{Biomateriales}
    \begin{itemize}[<+->]
      \item \textbf{Materiales:} Circonia y Alúmina de grado médico.
      \item \textbf{Propiedades Clave:} Biocompatibilidad y alta resistencia al
        desgaste.
      \item \textbf{Aplicaciones:} Implantes dentales y componentes de prótesis
        articulares (ej. cabezas femorales).
    \end{itemize}
  \end{block}
\end{frame}


\begin{frame}{Materiales para Sectores de Alto Rendimiento (2000 - Hoy)}
  La mejora en la tenacidad y fiabilidad ha permitido su uso en aplicaciones
  estructurales críticas.

  \begin{itemize}[<+->]
    \item \textbf{Aeroespacial:} Compuestos de Matriz Cerámica (CMC, ej. \ce{SiC/SiC})
      para componentes de motores a reacción, permitiendo mayor eficiencia de
      combustión.
      \vspace{0.5cm}
    \item \textbf{Blindaje y Protección:} Carburo de silicio (\ce{SiC}) y carburo
      de boro (\ce{B4C}) para blindaje personal y de vehículos por su extrema
      dureza y baja densidad.
      \vspace{0.5cm}
    \item \textbf{Generación de Energía:} Circonia estabilizada con itria (YSZ)
      como electrolito en celdas de combustible de óxido sólido (SOFC) por su
      conductividad iónica.
  \end{itemize}
\end{frame}

\begin{frame}{Cerámicas Superconductoras de Alta Temperatura (HTS)}
  Un descubrimiento disruptivo en la física de materiales fue la
  superconductividad en compuestos cerámicos.

  \pause

  \begin{block}{Descubrimiento Clave (1986)}
    Bednorz y Müller descubrieron superconductividad en un óxido cerámico de la
    familia de los cupratos a una temperatura superior a la de cualquier
    superconductor metálico conocido. Este hito les valió el \alert{Premio
    Nobel} de
    Física en 1987.
  \end{block}

  \pause

  \begin{block}{Material Principal y Propiedades}
    \begin{itemize}[<+->]
      \item \textbf{Ejemplo:} Óxido de Itrio, Bario y Cobre (YBCO), con fórmula
        \ce{YBa2Cu3O_{7-\delta}}.
      \item \textbf{Estructura:} Cristalina compleja de tipo perovskita.
      \item \textbf{Desafío Técnico:} Su naturaleza cerámica les confiere una
        alta fragilidad, dificultando su fabricación en forma de cables
        flexibles.
    \end{itemize}
  \end{block}
\end{frame}

\begin{frame}
  \begin{block}{Impacto y Aplicaciones}
    La capacidad de ser superconductores a temperaturas superiores a la del
    nitrógeno líquido (77 K, -196 °C) redujo drásticamente los costes de
    enfriamiento. Sus aplicaciones incluyen:
    \begin{itemize}[<+->]
        \item Imanes de alto campo para resonancia magnética (RM) y
          aceleradores de partículas.
        \item Limitadores de corriente de falla en redes eléctricas.
        \item Transporte por levitación magnética (Maglev).
     \end{itemize}
  \end{block}
\end{frame}

\begin{frame}{Fronteras Actuales: Computación Cuántica y Óptica Avanzada}
  \begin{block}{Computación Cuántica}
    \begin{itemize}[<+->]
      \item \textbf{Materiales:} Sustratos monocristalinos de Zafiro (\ce{Al2O3})
        y Nitruro de Silicio (\ce{Si3N4}).
      \item \textbf{Función:} Sus bajas pérdidas dieléctricas a temperaturas
        criogénicas son esenciales para fabricar resonadores que interactúan
        con cúbits superconductores, minimizando la decoherencia.
    \end{itemize}
  \end{block}

  \begin{block}{Cerámicas Transparentes}
    \begin{itemize}[<+->]
      \item \textbf{Materiales:} Oxinitruro de Aluminio (ALON), Espinela
        (\ce{MgAl2O4}).
      \item \textbf{Aplicación:} Ventanas para sensores militares y blindaje
        transparente, combinando dureza con transparencia óptica.
    \end{itemize}
  \end{block}
\end{frame}
